% ECONOMICS_PROBLEM: {"id": "Q58-103", "title": "Clean-technology policy sequences", "jel_code": "Q58", "source_id": "OP-0439"}
% !TeX program = xelatex
\documentclass[11pt,a4paper]{article}
\usepackage[margin=23mm]{geometry}
\usepackage{fontspec}
\setmainfont{Latin Modern Roman}
\usepackage{amsmath,amssymb,mathtools}
\usepackage{xcolor,enumitem,booktabs,longtable,array,xurl,hyperref}
\definecolor{muted}{HTML}{53636C}
\hypersetup{colorlinks=true,urlcolor=blue,linkcolor=blue}
\setlength{\parindent}{0pt}
\setlength{\parskip}{5pt}
\newcommand{\R}{\mathbb R}
\newcommand{\E}{\mathbb E}
\newcommand{\N}{\mathbb N}
\newcommand{\ind}{\mathbf 1}
\newcommand{\Var}{\operatorname{Var}}
\newcommand{\Cov}{\operatorname{Cov}}
\newcommand{\supp}{\operatorname{supp}}
\DeclareMathOperator*{\argmin}{arg\,min}
\DeclareMathOperator*{\argmax}{arg\,max}
\newcommand{\entrymeta}[3]{{\sffamily\small\color{muted}\textbf{\texttt{#1}}\quad|\quad #2\par #3\par}}
\newcommand{\Problemref}[1]{\texttt{#1}}
\newcommand{\JELref}[2]{\texttt{#2} (JEL \texttt{#1})}
\begin{document}
\section*{Clean-technology policy sequences}
\textbf{Problem ID:} \texttt{Q58-103}\quad \textbf{JEL:} \texttt{Q58}\par
% BEGIN ECONOMICS STATEMENT
\label{problem:OP-0439}
{\sffamily\small\color{muted}Original subject group: Climate, environment and energy\par}
\label{definitions:problem:30007721}
\textbf{Audit classification:} Background; proposed extension. Specified nonanticipating policies and expectation constraints for the stochastic optimization.
\subsection*{Definitions and precise research target}
\textit{Editorial formalization.} Consider an economy with fossil and clean production sectors over annual periods \(t=0,\ldots,30\). The state includes installed capital \(K_t\), clean-sector knowledge \(Z_t\), household assets, and political support \(S_t\). A policy path \(\pi\) specifies carbon taxes, clean-research subsidies, infrastructure investment, and transfers each year. Knowledge spillovers and irreversible investment are explicitly modeled. Let \(W(\pi)\) be discounted household welfare with stated distributional weights, and \(G_t(\pi)\) emissions. For cumulative emissions budget \(B\), fiscal limits \(F_t\), and implementability threshold \(s_0\), the proposed problem is
\[\pi^*\in\arg\max_{\pi\in\mathcal P}W(\pi)\quad\text{subject to}\quad E[\sum_{t=0}^{30}G_t(\pi)]\leq B,\quad E[\operatorname{Cost}_t(\pi)]\leq F_t,\quad E[S_t(\pi)]\geq s_0.\]
Here \(\mathcal P\) is the stated set of administratively feasible sequences, and \(\operatorname{Cost}_t\) is net public expenditure. Establish feasibility and whether the welfare supremum is attained; if it is not, report the supremum and approximating paths rather than assert an optimizer exists. Determine whether early innovation and infrastructure spending changes the welfare ranking of later carbon-tax paths after equilibrium entry and distributional effects. Estimate or bound learning, spillover, and support parameters instead of choosing them to justify a preferred sequence.

The state \(x_t=(K_t,Z_t,\text{assets}_t,S_t)\) follows a declared law \(x_{t+1}=f(x_t,u_t,\xi_{t+1})\), where controls \(u_t\) specify each tax, subsidy, public investment, and transfer. A policy is a nonanticipating measurable rule \(u_t=\pi_t(x_0,\xi_1,\ldots,\xi_t)\), rather than a sequence allowed to observe future shocks. With shocks present, use the displayed expected constraints, \(E\sum_tG_t\leq B\), \(E\operatorname{Cost}_t\leq F_t\), and \(E S_t\geq s_0\), with almost-sure constraints as a stricter alternative. Welfare \(W=E\sum_i\omega_i\sum_t\beta^tu_i(c_{it})\) uses fixed positive weights and integrable utilities. Political support is a defined model state or measured probability, with an explicit transition law. Optimization in a completely known finite menu is routine; parameter recovery and implementability motivate the agenda. All expectations use a stated baseline population and probability law. Identification means every admissible data-generating model with the same observed-data law yields the same displayed target; otherwise report the identified set. Exact equations, horizons and assumptions are editorial, rather than source-given canonical conjectures.
\subsection*{Variants and assumption changes}
\begin{enumerate}
\item \textbf{Editorial commitment variant.} Compare a committed thirty-year path with annually revised nonanticipating rules under a stated political transition model. Investor expectations must change consistently.
\item \textbf{Editorial risk-constraint variant.} Replace expectation constraints by \(P(\sum_tG_t\leq B)\geq1-\alpha\) and \(P(S_t\geq s_0\ \forall t)\geq1-\gamma\). The stricter feasible set may be empty. Fix \(\alpha,\gamma\in(0,1)\) before comparison.
\end{enumerate}
\subsection*{References and source audit}
\textbf{1.} Broad transformation and policy discussion; precise sequence optimization is editorial. Locator: Economic Journal 132(644), 1259–1289; author-institution publication summary. Publisher retrieval failed; author-institution metadata and summary inspected. URL: \url{https://academic.oup.com/ej/article/132/644/1259/6519262}.

\textbf{2.} Climate macroeconomics survey; no exact fifty-year or tipping-event conjecture verified. Locator: NBER working-paper record and abstract; AEA forthcoming record, DOI 10.1257/jel.20261831. Primary indexed record verified; direct NBER/PDF retrieval failed. AEA forthcoming abstract inspected; no final journal publication year asserted. URL: \url{https://www.nber.org/papers/w33567}.
\begin{enumerate}\raggedright
\item\hypertarget{definitions:source:30007721:1}{} Nicholas Stern. 2022. \emph{A Time for Action on Climate Change and a Time for Change in Economics}. Economic Journal 132(644), 1259–1289; author-institution publication summary.
\url{https://academic.oup.com/ej/article/132/644/1259/6519262}.
DOI: \url{https://doi.org/10.1093/ej/ueac005}.
\item\hypertarget{definitions:source:30007721:2}{} Adrien Bilal; James H. Stock. 2025. \emph{A Guide to Macroeconomics and Climate Change}. NBER working-paper record and abstract; AEA forthcoming record, DOI 10.1257/jel.20261831.
\url{https://www.nber.org/papers/w33567}.
DOI: \url{https://doi.org/10.3386/w33567}.
\end{enumerate}

\subsection*{Common conventions: Source mathematical conventions}

These conventions apply to each entry unless explicitly replaced.

\textbf{Models and identification.} A primitive model \(m\) specifies all probability laws, feasible choices, information, equilibrium and measurement rules used by its target. The admissible class \(\mathcal M\) is fixed independently of outcomes. If \(P_O(m)\) is its observable probability law and \(\tau(m)\) the target, its identified set at observable law \(P\) is
\[
 \mathcal I_\tau(P)=\{\tau(m):m\in\mathcal M,\ P_O(m)=P\}.
\]
Point identification means that this set is a singleton. An empty set signals incompatibility of data and assumptions. Coordinate bounds need not describe the joint set. Bounds are sharp only if every retained target is attainable or arbitrarily approachable by compatible models. Finite bounds require explicit restrictions on unobserved quantities; unrestricted confounding, hidden wealth, extrapolation or arbitrary equilibrium selection can yield uninformative sets.

\textbf{Causal interventions.} A potential outcome \(Y_i(p)\) is the outcome under a complete policy regime \(p\), including timing, financing, information and market spillovers. The target population and units are fixed before treatment. With interference, use the complete assignment vector or a justified exposure mapping; individual no-interference notation alone cannot identify an economy-wide intervention. A difference of mean potential outcomes does not require the cross-world joint distribution. Conditioning on post-treatment migration, survival, enrollment or employment changes the estimand and needs separate justification.

\textbf{Statistical statements.} An estimator, confidence set, forecast or test specifies a sampling law, model class and loss/coverage criterion. Uniform \((1-\alpha)\)-coverage means \(\inf_{m\in\mathcal M}\Pr_m\{\tau(m)\in C_n\}\geq1-\alpha\), or an explicitly asymptotic version. The whole parameter space trivially has coverage, so informativeness requires a stated width, risk or power objective. A forecasting improvement is relative to a named benchmark, information set and evaluation population.

\textbf{Equilibrium and optimization.} An equilibrium comprises feasible plans, optimality, beliefs where needed, budget constraints and all market/resource clearing. The primitives or a stated benchmark define each correspondence \(\mathcal E(p;m)\). Nonemptiness is assumed or proved, never inferred just from compact policy sets. A selected-equilibrium target uses a declared selection rule; a robust target quantifies over all permitted equilibria. Use suprema or infima unless attainment is established. An \(\varepsilon\)-optimal policy is feasible and within \(\varepsilon\) of the finite optimum value. Report an empty feasible set separately.

\textbf{Welfare and accounting.} Normative utility, interpersonal normalization, social weights, numeraire, discounting and the use of public receipts are primitives. Behavioral utility need not coincide with the welfare criterion. Household welfare includes ownership income and rents once; adding profits again to compensating variation generally double-counts them. A monetary equivalent is defined by an explicit transfer and price convention, with monotonicity and range conditions. Average effects, mechanism contributions, transfers and welfare losses are different targets. Infinite sums require convergence and dynamic feasible plans require terminal or no-Ponzi conditions.

\textbf{Source versions.} A working-paper date, revision date and journal publication year are distinguished. A DOI for a journal version is not automatically the DOI of an earlier manuscript. Locators refer to the stated version and printed pages where specified. A metadata-only check does not verify a claim about a full-text conclusion. Every entry's source subsection narrows the provenance claim accordingly.
% END ECONOMICS STATEMENT
\end{document}
